% !TeX root = ../informe.tex
% ===========================================================================
%  3. Alcance funcional
%
%  Qué hace el sistema y qué deliberadamente no hace. Es contexto para las
%  secciones de arquitectura, no el capítulo principal: basta con recoger
%  las funcionalidades de forma compacta y señalar lo que se acotó.
% ===========================================================================

\section{Alcance funcional}
\label{sec:alcance}

\subsection{Roles y permisos}
\label{sec:roles}

\instruccion{Una matriz de funcionalidad por rol se lee de un vistazo y
evita párrafos repetitivos. Marca con «Sí» y «No»; si algún permiso depende
de una condición (por ejemplo, actuar solo sobre lo propio), dilo en la
celda. Añade una columna por rol que exista de verdad.}

\begin{table}[H]
  \centering
  \caption{Matriz de roles y permisos.}
  \label{tab:roles}
  \begin{tabular}{L{6.2cm} C{2.6cm} C{2.6cm}}
    \toprule
    \textbf{Funcionalidad} & \ph{Rol 1} & \ph{Rol 2} \\
    \midrule
    \ph{...} & \ph{...} & \ph{...} \\
    \bottomrule
  \end{tabular}
\end{table}

\subsection{Módulos y pantallas}
\label{sec:modulos}

\instruccion{Los módulos funcionales con sus pantallas. Conviene que esta
tabla case con las piezas de la \cref{tab:componentes}: si un módulo no
corresponde a ninguna pieza, o al revés, es señal de que la descomposición
no está clara.}

\begin{table}[H]
  \centering
  \caption{Módulos y pantallas del sistema.}
  \label{tab:modulos}
  \begin{tabular}{L{3.2cm} L{4.4cm} L{5.4cm}}
    \toprule
    \textbf{Módulo} & \textbf{Pantalla} & \textbf{Descripción} \\
    \midrule
    \ph{...} & \ph{...} & \ph{...} \\
    \bottomrule
  \end{tabular}
\end{table}

\subsection{Fuera de alcance}
\label{sec:fuera-alcance}

\instruccion{Enumera lo excluido. No es una disculpa: delimitar el alcance
es una decisión de diseño, y declararla evita que se lea como una carencia.}

\ph{Funcionalidades excluidas y el motivo.}
